\section{Introduction}
\label{sec:intro}

Most industrial robots use simple end-effectors, such as parallel grippers, because they are reliable, inexpensive, and easy to control. They suffice for pick-and-place, but once an object is grasped they offer little control over its pose within the grasp. That pose often matters. A part picked from a bin wherever it was easiest to grasp must frequently be presented to a fixture, a machine, or an assembly in one particular grasp; today this requires a regrasp station, a dedicated jig, or a more dexterous hand. A gripper that could adjust its grasp in place would remove that step, and it would also relax the constraints on the initial grasp.

People face the same limitation with one hand and overcome it with a second hand, even one that does very little. Holding an object still while the fingers slide along it, or giving it a point to pivot about, is often enough. A passive tool in the workspace can play this role: contact with the tool holds or guides the object while the robot moves the gripper. This makes the workspace a design variable for in-hand manipulation. The question is then what geometry to add, where to place it, and how to move the robot so that a desired change of grasp follows from those contacts. Our thesis is that this question has a closed-form answer. Instead of giving a simple gripper more sensing, actuation, or planning, we give its workspace the geometry it is missing, and we show that this geometry can be designed in closed form.

External assistance is an established route to dexterity with simple grippers. Extrinsic dexterity uses forces from gravity, inertia, and environmental contacts to change an object's pose within the grasp \cite{dafle2014extrinsic}. Prehensile pushing constructs such changes by pressing a grasped object against fixed features in the environment \cite{chavan2015prehensile,chavan2018stable}; sampling-based planners and motion cones chain these pushes into regrasps, including motions that combine sliding and turning, by modelling the frictional contacts between gripper, object, and pusher \cite{chavan2017sampling,chavan2020planar}. Fixtureless fixturing pushes the object against an environmental feature that holds it still while the fingers slide over it \cite{chavan2018fixtureless}, the closest precedent for the shift studied here. Pivoting against the environment and shared grasping reorient objects through contact with a support \cite{hou2019reorienting,hou2020manipulation}, inertial loads can drive in-grasp sliding \cite{shi2017dynamic}, and visual and tactile feedback can regulate pivoting \cite{karayiannidis2016adaptive,toskov2023hand,kim2024texterity}. These methods establish that interaction with the environment extends a gripper's capabilities. They take the environment as given, however, and obtain the motion from friction coefficients, grip-force values, numerical planning, or feedback (\tabref{tab:positioning}).

\label{sec:related}
Designing the constraints themselves has a long history. Passive fences orient parts, modular fixtures constrain them through deliberately chosen contact geometry, and sensorless orienting uses geometric contact sequences \cite{wiegley1997complete,brost1996complete,goldberg1993orienting,asada1985fixturing}; their objective is a final part pose, not a change of pose while the part remains in a gripper. Effector form design derives a finger's shape from the motion it must produce \cite{rodriguez2013effector}, and bearings or passive rollers in the fingers make pivoting and sliding within the grasp mechanically easier \cite{carlisle1994pivoting,fisk2026rollers}; these approaches modify the gripper, whereas we leave it unchanged and design a fixture in the workspace. Open-loop pivoting shows that a simple gripper can reorient an object without feedback \cite{holladay2015general,cruciani2017hand}. Our earlier assistive tool used two planes with their angle selected through simulation \cite{bolduc2023inhand}; here the tool's geometry and placement follow from analysis instead. The analysis draws on classical planar contact mechanics: instantaneous rotation centres on contact normals \cite{mason2001mechanics}, first- and second-order contact constraints \cite{ohwovoriole1981extension,rimon1998mobility}, the mechanics of pushing \cite{mason1986mechanics,peshkin1988planning,lynch1996stable}, and limit surfaces for sliding contact patches \cite{goyal1991planar,howe1996practical,xydas1999modeling}. These ingredients are classical; what has been missing is their use as a design method for workspace geometry. Among the approaches compared in \tabref{tab:positioning}, ours is the only one that designs the workspace and changes the grasp in place, and the only one in which no friction coefficient, grip-force value, numerical planner, or feedback enters the manipulation motion.

\begin{table}[htbp]
    \centering
    \footnotesize
    \renewcommand{\arraystretch}{1.15}
    \setlength{\tabcolsep}{4pt}
    \begin{tabular}{>{\raggedright\arraybackslash}p{3.3cm}*{5}{>{\centering\arraybackslash}p{1.62cm}}>{\centering\arraybackslash}p{2.0cm}}
        \hline
        Approach & Same gripper & Workspace designed & Grasp changed in place & Friction or force in the plan & Feedback in motion & Motion computed by \\
        \hline
        Fences, modular fixtures, sensorless orienting \cite{wiegley1997complete,brost1996complete,goldberg1993orienting} & yes & yes & no & -- & no & complete search \\
        Prehensile pushing, motion cones, fixtureless fixturing \cite{chavan2015prehensile,chavan2017sampling,chavan2018fixtureless,chavan2020planar} & yes & no & yes & yes & no & sampling, optimization \\
        Pivoting, shared grasping \cite{hou2019reorienting,hou2020manipulation} & yes & no & yes & yes & force control & contact-mode planning \\
        Feedback-regulated pivoting \cite{karayiannidis2016adaptive,toskov2023hand,kim2024texterity} & yes & no & yes & estimated & vision, touch & controller \\
        Rollers or bearings in the fingers \cite{carlisle1994pivoting,fisk2026rollers} & no & no & yes & varies & varies & varies \\
        \textbf{Two free rods (this work)} & \textbf{yes} & \textbf{yes} & \textbf{yes} & \textbf{no} & \textbf{no} & \textbf{closed form} \\
        \hline
    \end{tabular}
    \caption{What approaches to changing a grasp with a simple gripper require. ``Friction or force in the plan'': whether a friction coefficient or a grip-force value is needed to compute the motion. ``Feedback in motion'': sensing used while the object is being manipulated, not to locate it beforehand.}
    \label{tab:positioning}
\end{table}

The proposed tool consists of two rods that spin freely on bearings (\figref{fig:tworod}). The gripper pinches a planar object across its thickness and first brings it against both rods. During a shift, the rods hold the object still while the gripper slides along it. During a turn, the gripper orbits one rod and the object pivots within the grasp. The wrist orientation stays fixed during each manipulation motion. A requested planar displacement is decomposed into a shift followed by a turn; when both rods contact the same edge, the object is carried away and re-engaged before turning. The tool amounts to two pins on bearings: it leaves the gripper and its controller unchanged, and the manipulation motions themselves need no sensing.

The paper makes four contributions:
\begin{itemize}
    \item \emph{A design principle for passive in-grasp manipulation.} We show why passive contacts with a fixed wrist cannot uniquely determine simultaneous sliding and turning within the grasp. This motivates separating a planar change of grasp into a shift and a turn (Section~\ref{sec:design}).
    \item \emph{A minimal tool with a closed-form design.} We establish that two rods are the minimum needed to support an arbitrary shift in one motion, and show how the same pair supports a subsequent turn. We derive the rod placements, gripper paths, and conditions under which an installed pair can perform a requested change of grasp (Sections~\ref{sec:method}, \ref{sec:tworod}, and~\ref{sec:installation}).
    \item \emph{Motion computation without friction or force parameters.} The manipulation motions require neither friction coefficients nor grip-force values as planning inputs. Freely rotating rods also allow grasp slip during a turn to be predicted from the finger-pad size and the lever arm (Section~\ref{sec:turn}).
    \item \emph{Experimental and simulation evidence.} We validate turns on hardware and evaluate shifts and complete shift--turn sequences in simulation, comparing the results with the analytical predictions (Section~\ref{sec:experiments}).
\end{itemize}

Execution assumes a known initial grasp pose and grip settings that retain the object while permitting the required sliding and pivoting; contact, grasp travel, clearance, and robot reach must also be checked. Hardware trials with a single rod demonstrate the turn in the vertex-contact case covered by the analysis, and shifts and complete sequences are tested in simulation.

\begin{figure}[htbp]
\centering
% Computed from the shift and turn constructions for a convex pentagon (mm, 1 mm = 0.05 cm):
% Delta = 80 deg, d = (6,18) mm, R = 5 mm, lambda = 34 mm. Written by python -m geometry; do not edit coordinates by hand.
\begin{tikzpicture}[paperfigure,
  gripbody/.style={fill=black!55,draw=black!80,line width=.4pt},
  grippad/.style={fill=black!75,draw=black},
  ghostobj/.style={draw=objectblue!60,densely dashed,fill=none,line width=.5pt}]
\newcommand{\gripper}[3]{%
  \begin{scope}[shift={(#1,#2)},#3]
    \fill[gripbody,rotate=30] (-3,0) rectangle (3,-24);
    \fill[grippad] (0,0) circle[radius=4];
    \fill[white] (0,0) circle[radius=1];
  \end{scope}}
\path[use as bounding box] (0,-0.10) rectangle (15.6,11.1);

\node[paneltitle] at (0,10.9) {(a) Engage};
\begin{scope}[shift={(4.084,9.114)},x=0.05cm,y=0.05cm]
\draw[ghostobj] (-28.26,-21.01) -- (35.74,-25.01) -- (51.74,-2.01) -- (23.74,17.99) -- (-34.26,8.99) -- cycle;
\draw[object] (-34.00,-17.00) -- (30.00,-21.00) -- (46.00,2.00) -- (18.00,22.00) -- (-40.00,13.00) -- cycle;
  \draw[feature] (-35.76,18.72) circle[radius=5]; \fill (-35.76,18.72) circle(.6pt);
  \draw[feature] (-43.12,3.09) circle[radius=5]; \fill (-43.12,3.09) circle(.6pt);
  \draw[motion,line width=.9pt] (30.00,-26.00) -- (20.16,-19.13);
  \gripper{0.00}{0.00}{}
  \node[anchor=south] at (-35.76,24.22) {$a$};
  \node[anchor=east] at (-48.62,3.09) {$b$};
\end{scope}
\node[paneltitle] at (8.0,10.9) {(b) Shift};
\begin{scope}[shift={(11.898,9.071)},x=0.05cm,y=0.05cm]
\draw[object] (-34.00,-17.00) -- (30.00,-21.00) -- (46.00,2.00) -- (18.00,22.00) -- (-40.00,13.00) -- cycle;
  \draw[guide] (-39.57,5.87) -- (25.72,-3.81);
  \node[anchor=south west,inner sep=1pt,text=black!60] at (11.72,-2.31) {grasp line};
  \draw[feature] (-35.76,18.72) circle[radius=5]; \fill (-35.76,18.72) circle(.6pt);
  \draw[feature] (-43.12,3.09) circle[radius=5]; \fill (-43.12,3.09) circle(.6pt);
  \draw[guide] (-35.00,13.78) -- (-33.63,4.99);
  \draw[force] (-37.45,29.59) -- (-35.07,14.27);
  \draw[guide] (-38.21,4.07) -- (-33.63,4.99);
  \draw[force] (-53.90,0.93) -- (-38.70,3.97);
  \fill (-33.63,4.99) circle(1.2pt); \node[anchor=north west,inner sep=1pt] at (-32.83,4.19) {$\mathbf x$};
  \gripper{0.00}{0.00}{,opacity=.28}
  \draw[motion,line width=1pt] (-2.97,0.44) -- (-15.80,2.34);
  \node[anchor=south,inner sep=2pt] at (-9.38,3.39) {$-\mathbf d_0$};
  \gripper{-18.77}{2.78}{}
\end{scope}
\node[paneltitle] at (0,6.764132811653769) {(c) Turn};
\begin{scope}[shift={(4.278,1.647)},x=0.05cm,y=0.05cm]
\draw[ghostobj] (-34.00,-17.00) -- (30.00,-21.00) -- (46.00,2.00) -- (18.00,22.00) -- (-40.00,13.00) -- cycle;
\draw[object,fill opacity=.45,draw opacity=.45] (-11.45,-7.51) -- (40.14,30.56) -- (37.62,58.47) -- (3.31,55.79) -- (-35.33,11.61) -- cycle;
\draw[object] (-0.28,14.25) -- (14.77,76.58) -- (-5.10,96.33) -- (-29.66,72.23) -- (-30.87,13.55) -- cycle;
  \draw[feature] (-35.76,18.72) circle[radius=5]; \fill (-35.76,18.72) circle(.6pt);
  \draw[feature] (-43.12,3.09) circle[radius=5]; \fill (-43.12,3.09) circle(.6pt);
  \draw[motion,line width=1pt] (-18.77,2.78) arc[start angle=-43.156,end angle=36.844,radius=23.296];
  \gripper{-18.77}{2.78}{,opacity=.28}
  \gripper{-12.50}{17.43}{,opacity=.28}
  \gripper{-17.12}{32.69}{}
  \node[anchor=east] at (-41.26,21.72) {$a$};
  \node[anchor=east] at (-48.62,3.09) {$b$};
\end{scope}
\node[paneltitle] at (8.0,6.764132811653769) {(d) Result};
\begin{scope}[shift={(11.550,3.282)},x=0.05cm,y=0.05cm]
\draw[ghostobj] (-34.00,-17.00) -- (30.00,-21.00) -- (46.00,2.00) -- (18.00,22.00) -- (-40.00,13.00) -- cycle;
\draw[object] (16.84,-18.44) -- (31.89,43.90) -- (12.02,63.65) -- (-12.54,39.55) -- (-13.75,-19.13) -- cycle;
  \fill[objectblue] (6.00,18.00) circle(1.3pt);
  \draw[motion,line width=.9pt] (0,0) -- (5.52,16.56);
  \node[anchor=east,inner sep=2pt] at (3.30,9.90) {$\mathbf d$};
  \draw[motion,line width=.9pt] (26,-4) arc[start angle=-10,end angle=70,radius=26];
  \node[anchor=west] at (31,14) {$\Delta=80^\circ$};
  \gripper{0.00}{0.00}{}
\end{scope}
\draw[black!15] (7.8,-0.10) -- (7.8,11.0); \draw[black!15] (0,7.01) -- (15.6,7.01);
\end{tikzpicture}
\caption{Changing the grasp with two rods, using the constructions of Sections~\ref{sec:method} and~\ref{sec:turn}. (a)~The gripper brings the plate against two installed rods that spin freely on bearings. Dark disk: gripper pad; bar: wrist, whose orientation stays fixed throughout. (b)~The rod forces (orange) act along push lines meeting at $\mathbf x$ on the grasp line: the plate stays still while the gripper slides by $-\mathbf d_0$. (c)~The gripper orbits rod~$a$; the plate pivots within the grip, turns by $\Delta$, and separates from rod~$b$. (d)~In the gripper frame, the plate has turned by $\Delta=80^\circ$ and shifted by the requested $\mathbf d=\mathbf Q(\Delta)\mathbf d_0$. Teal arrows: gripper motion in (a--c), relative object displacement and rotation in (d); dashed outlines: initial poses.}
\label{fig:tworod}
\end{figure}
